\section*{CHAPTER 1: INTRODUCTION}
\addcontentsline{toc}{section}{CHAPTER 1: INTRODUCTION}

\subsection*{1.1 Background of the Study}
The rapid acceleration of digital technology has created a socio-economic paradox in developing municipalities. While digital literacy is now a mandatory prerequisite for global and local employment, the infrastructure required to teach these skills is lagging. In many areas, the Technical Education and Skills Development Authority (TESDA) and local universities are under immense pressure to provide computer-based certifications, yet they are hamstrung by acute capital scarcity. Conversely, the widespread accessibility of smartphones has triggered an economic decline in the traditional internet café (PC bang) industry, leaving high-specification hardware idle during the day.

In the Municipality of Maigo, Lanao del Norte, this paradox is highly pronounced. The municipality hosts the Mindanao State University - Maigo College of Education, Science and Technology (MSU-MCEST) and local TESDA operations, both of which face significant infrastructural constraints. MSU-MCEST’s computer laboratories primarily cater to basic introductory courses, lacking the capacity for advanced, certification-level training. Similarly, local TESDA programs cannot offer crucial ICT certifications because they lack the specialized laboratories required for training and assessment.

Simultaneously, Maigo’s internet café owners are facing bankruptcy. As their traditional demographic—students seeking basic research or entertainment—migrates to mobile data, their high-performance hardware sits dormant during standard business hours (8:00 AM to 5:00 PM), accruing depreciation without generating revenue.

This study posits that these isolated struggles represent a significant opportunity for Collaborative Governance. By institutionalizing a Public-Private Partnership (PPP) that transforms struggling internet cafés into government-contracted, ``Dual-Use Tech Spaces,'' the LGU can bridge the digital divide through a ``Zero-CapEx'' model while providing an economic lifeline to local MSMEs.

\subsection*{1.2 Statement of the Problem}
This study aims to develop a localized public policy framework to institutionalize the dual-use of commercial internet cafés as government-contracted training micro-laboratories during the day and commercial hubs at night in the Municipality of Maigo.

Specifically, it seeks to answer the following:
\begin{enumerate}
    \item \textbf{The Governance and Legal Factor:} What existing local ordinances, zoning laws, or procurement guidelines in Maigo currently hinder or support the LGU in formally leasing private IT infrastructure?
    \item \textbf{The Business and Economic Factor:} What financial models (e.g., service contracting, utility subsidies) are required to make government leasing a viable economic lifeline for Maigo’s MSMEs?
    \item \textbf{The Infrastructure and Technical Factor:} How do the existing hardware and network architectures (such as Diskless/PXE-Boot systems) in local cafés measure against the baseline standards mandated by TESDA and MSU-MCEST?
    \item \textbf{The Policy Output:} Based on the findings, what comprehensive LGU ordinance and implementation plan can be proposed to operationalize the ``Dual-Use Tech Space'' program?
\end{enumerate}

\subsection*{1.3 Objectives of the Study}
\textbf{General Objective:} To design a formalized PPP framework and municipal ordinance that utilizes private internet cafés in Maigo as official TESDA and university-accredited computer laboratories.

\textbf{Specific Objectives:}
\begin{itemize}
    \item To assess the current IT infrastructure gap within MSU-MCEST and TESDA centers in Maigo.
    \item To evaluate the technical readiness and economic viability of local internet cafés to serve as formal educational spaces.
    \item To identify the legal and bureaucratic mechanisms under the Philippine PPP Code (RA 11966) for daytime leasing agreements.
    \item To draft a localized policy recommendation that addresses the digital divide through collaborative asset utilization.
\end{itemize}

\subsection*{1.4 Significance of the Study}
\begin{itemize}
    \item \textbf{The Local Government of Maigo:} Provides a Zero-Capital-Expenditure (Zero-CapEx) model to upgrade local digital infrastructure without building redundant facilities.
    \item \textbf{Internet Café Owners (MSMEs):} Offers a sustainable, institutionalized revenue stream to maximize hardware ROI and prevent business closure.
    \item \textbf{TESDA and MSU-MCEST:} Grants immediate access to high-spec, maintained laboratories, expanding their capacity to offer advanced ICT certifications.
    \item \textbf{The Citizens of Maigo:} Democratizes access to digital skills training, directly combating the digital divide and improving employability for out-of-school youth and workers.
\end{itemize}

\subsection*{1.5 Scope and Delimitations}
This study focuses strictly on the public administration and policy formulation aspects of utilizing private internet cafés for public education in Maigo. The subjects include registered café owners, representatives from MSU-MCEST, TESDA administrators, and LGU officials. The research will not assess pedagogical effectiveness or conduct the actual training; it focuses exclusively on the governance mechanism of providing the space.

\subsection*{1.6 Conceptual Framework}
The study is anchored on Collaborative Governance and New Public Management (NPM). The LGU acts as the ``Central Orchestrator,'' using its legislative power to create accreditation systems and its executive power to facilitate service contracts. This bridges the ``Educational Deficit'' (lack of hardware) with the ``Private Sector Surplus'' (underutilized hardware), resulting in an institutionalized policy that serves the public good while sustaining private enterprise.